\chapter{Mar.~18 --- KZ Equations, Part 2}

\section{Correlation Functions, Continued}

\begin{prop}
  The correlation functions span the
  space of solutions of the KZ equations.
\end{prop}

\begin{remark}
  Let $\xi = (\lambda_0, \dots, \lambda_N; g_1, \dots, g_N; u_0)$,
  $\lambda_i \in \mathfrak{h}^*$,
  $g_i \in \Hom_{\g}(L_{\lambda_i}, L_{\lambda_{i - 1}} \otimes V_{\mu_i})$
  where $u_0 \in L_{\lambda_0}^*$. Let
  \[
    \Xi
    = \bigoplus_{\lambda_0, \dots, \lambda_{N - 1}} \bigg(
      L_{\lambda_0}^*
      \otimes \bigotimes_{i = 1}^N \Hom_{\g}
      (L_{\lambda_i}, L_{\lambda_{i - 1}} \otimes V_{\mu_i})
    \bigg).
  \]
  The only nonzero terms satisfy
  $\lambda_i + \mu_i - \lambda_{i - 1} \in Q$,
  and this space has a natural
  $\g$-module structure, and
  the $\Hom_{\g}(L_{\lambda_i}, L_{\lambda_{i - 1}} \otimes V_{\mu_i})$ are
  trivial $\g$-modules. Each $\xi \in \Xi$
  defines a unique correlation function
  \[
    \psi^\xi(z_1, \dots, z_N)
    \in V
    = V_{\mu_1} \otimes \cdots \otimes V_{\mu_N} \otimes L_{\lambda_N}^*.
  \]
\end{remark}

\begin{theorem}\label{thm:correlation-functions-span-solutions}
  Fix $\lambda_N$ and
  $\mu_i$ for $i = 1, \dots, N$. Let
  $V_{\mu_i}$ be the lowest weight modules
  satisfying either of the two
  following conditions:
  \begin{enumerate}
    \item \emph{Generic case:}
      The weights $\lambda_N, \lambda_N + \mu_N, \dots, \lambda_N + \mu_N + \cdots + \mu_1$
      are generic and $k$ is generic with
      respect to each;
    \item \emph{Finite-dimensional case:}
      $\lambda_N, \mu_i \in P^+$,
      $k \notin \Q$, and the $V_{\mu_i}$
      are irreducible finite-dimensional
      modules with lowest weight $-\mu_i$.
  \end{enumerate}
  Then the correlation functions span
  the space of $V$-valued solutions
  of the KZ equations in a region $D$.
  Moreover, the following map
  \begin{align*}
    \Xi &\longrightarrow \Gamma(D, V_{\mathrm{KZ}}) \\
    \xi &\longmapsto \psi^\xi(z_1, \dots, z_N)
  \end{align*}
  is an isomorphism of $\g$-modules.
\end{theorem}

\begin{remark}
  In Case 1 of Theorem \ref{thm:correlation-functions-span-solutions},
  the $\lambda_i$ are generic, and in
  Case 2, the $\lambda_i \in P^+$.
\end{remark}

\begin{lemma}
  Under assumptions (1) or (2) of
  Theorem \ref{thm:correlation-functions-span-solutions},
  the map
  \[
    (\lambda_0, \dots, \lambda_{N - 1}; g_1, \dots, g_N; u_0)
    \longmapsto
    \langle u_0, g_1 \cdots g_N \cdot \rangle
  \]
  is an isomorphism of $\g$-modules.
\end{lemma}

\begin{proof}
  Note that $\Hom_\g(L_{\lambda_i}, L_{\lambda_{i - 1}} \otimes V_{\mu_i}) \cong \Hom_\g(L_{\lambda_{i - 1}}^*, V_{\mu_i} \otimes L_{\lambda_i}^*)$.
  So it suffices to show that
  \[
    \bigoplus_v \Hom_\g(L_v^*, V_\mu \otimes L_\lambda^*) \otimes L_v^*
    \longrightarrow
    V_\mu \otimes L_\lambda^*
  \]
  is an isomorphism. This is true if
  we are in the finite-dimensional case.
  In general, if $\lambda$ and $\lambda + \mu$
  are generic, the $L_v^*$ being
  irreducible implies that the above map
  is injective. Since $v$ is
  generic, $L_v^* \cong M_v^*$
  and $L_\lambda^* \cong M_\lambda^*$, so
  $\Hom_\g(L_v^*, V_\mu \otimes L_\lambda^*) \cong V_\mu^{(\lambda - v)}$
  (the weight subspace).
\end{proof}

\begin{corollary}
  If $u_0$ is the lowest weight vector, then
  the correlation functions span solutions
  with values in $V^{\mathfrak{n}_-}$.
  Moreover, if $u_0$ is $\g$-invariant,
  then the correlation functions take
  values in $V^\g$.
\end{corollary}

\begin{remark}
  We had an identification
  $\Gamma(D, V_{\mathrm{KZ}}) \cong V$
  via $\psi \mapsto \psi(p)$
  (i.e. evaluation at $p$). Now we
  have
  \[
    \Gamma(D, V_{\mathrm{KZ}})
    \cong \Xi \cong V,
  \]
  which is a more sophisticated
  identification that is not attached to
  a particular point.
\end{remark}

\section{Trigonometric KZ Equations}

\begin{remark}
  Let $\phi(z_1, \dots, z_N) = \langle u_0, \psi(z_1, \dots, z_N) u_{N + 1} \rangle \in V_{\mu_1} \otimes \cdots \otimes V_{\mu_N}$.
  Then
  \[\phi \in (V_{\mu_1} \otimes \cdots \otimes V_{\mu_N})^{\lambda_N - \lambda_0},\]
  and $\phi$ satisfies the following
  differential equation
  \[
    (k + h^\vee) \frac{\partial}{\partial z_i} \phi
    = \bigg(
      \sum_{j = 1}^N \frac{\Omega_{ij}}{z_i - z_j}
    \bigg) \phi
    - \frac{1}{z_i} \sum_e (a_e)_i \langle u_0, \psi(z_1, \dots, z_N) a^\ell u_{N + 1} \rangle,
  \]
  where $a^e, a_e$ are dual basis
  elements in $\g$. For
  $\alpha \in R_+$, let
  $e_\alpha \in \g^\alpha$,
  $f_\alpha \in \g^{-\alpha}$ such that
  $\langle e_\alpha, f_\alpha \rangle = 1$. Then
  \[
    \Omega = \sum_{\alpha \in R_+}
    (e_\alpha \otimes f_\alpha + f_\alpha \otimes e_\alpha)
    + \sum_p x_p \otimes x_p,
  \]
  where $x_p$ is the dual
  basis in $\mathfrak{h}$, and
  $C = \sum_{\alpha \in R_+} (e_\alpha f_\alpha + f_\alpha e_\alpha) + \sum_p x_p^2$.
  We can rewrite
  \begin{align*}
    \sum_e (a_e)_i \langle u_0, \psi a^\ell u_{N + 1} \rangle
    &= \sum_{\alpha \in R_+} (e_\alpha)_i
    \langle u_0, \psi f_\alpha u_{N + 1} \rangle
    + \sum_{p} (x_p)_i \rangle u_0, \psi \langle x_p, \lambda_N \rangle u_{N + 1} \rangle \\
    &= \sum_{\alpha \in R_+} (e_\alpha)_i
    \sum_{j = 1}^N (f_\alpha)_j \phi + (h_{\lambda_N})_i \phi.
  \end{align*}
  Note that $u_0$ is $\mathfrak{n}_-$-invariant.
  Now $\phi$ has weight $\lambda_N - \lambda_0$,
  so
  \[
    \sum_p \sum_{j = 1}^N (x_p)_i (x_p)_j \phi
    = (h_{\lambda_N - \lambda_0})_i \phi.
  \]
  Then we can rewrite
  \pagebreak
  \begin{align*}
    \sum_e (a_e)_i \langle u_0, \psi a^\ell u_{N + 1} \rangle
    &= \sum_{j = 1}^N
    \bigg(
      \sum_{\alpha \in R_+} (e_\alpha)_i (f_\alpha)_j
      + \frac{1}{2} \sum_p (x_p)_i (x_p)_j
    \bigg) \phi
    + \frac{1}{2} (h_{\lambda_N} + \lambda_0)_i \phi \\
    &= \sum_{j = 1}^N \Omega_{i, j}^+ \phi
    + \frac{1}{2}(h_{\lambda_N + \lambda_0})_i \phi
    + \bigg(\sum_{\alpha \in R_+} e_\alpha f_\alpha + \frac{1}{2} \sum_p x_p^2\bigg)_i \phi,
  \end{align*}
  where $\Omega^+ = \sum_{\alpha \in R_+} e_\alpha \otimes f_\alpha + \frac{1}{2} \sum_p x_p \otimes x_p$
  \[
    \Omega^- = \Omega - \Omega^+
    = \sum_{\alpha \in R_+} f_\alpha \otimes e_\alpha
    + \frac{1}{2} \sum_p x_p \otimes x_p.
  \]
  Check as an exercise that
  $[e_\alpha, f_\alpha] = h_\alpha$ and
  that
  \[
    \sum e_\alpha f_\alpha + \frac{1}{2} \sum_p x_p^2
    = \frac{1}{2} C + h_\rho.
  \]
  Then we have that
  \begin{align*}
    \sum_\ell (a_\ell)_i \langle u_0, \psi a^\ell u_{N + 1} \rangle
    &= \sum_{\substack{j = 1 \\ j \ne i}}^N \Omega_{i, j}^+ \phi
    + \frac{1}{2} (h_{\lambda_N + \lambda_0 + 2\rho})_i \phi
    + \frac{1}{2}(C)_i \phi \\
    &= \bigg(
      \sum_{j = 1}^N \Omega_{i, j}^+
      + \frac{1}{2}(h_{\lambda_N + \lambda_0 + 2\rho})_i
      + \frac{1}{2} \langle \mu_i, \mu_i + 2\rho \rangle
    \bigg) \phi.
  \end{align*}
  Thus we can rewrite the KZ equations as
  \[
    (k + h^\vee) z_i \partial_{z_i} \phi
    = \bigg(\sum_{\substack{j = 1 \\ j \ne i}}^N
    \frac{z_i \Omega_{i, j}^- + z_j \Omega_{i, j}^+}{z_i - z_j}
    - \frac{1}{2}(h_{\lambda_N + \lambda_0 + 2\rho})_i
    - \frac{1}{2}\langle \mu_i, \mu_i + 2\rho \rangle\bigg) \phi
  \]
\end{remark}

\begin{theorem}
  Let $z_i = e^{z_i}$. Then
  \[
    (k + h^\vee)
    \partial_{x_i} \phi
    = \bigg(
      \sum_{\substack{j = 1 \\ j \ne i}}^N
      r_{i, j}(x_i - x_j) -
      \frac{1}{2} (h_{\lambda_N + \lambda_0 + 2\rho})_i
      - \langle \mu_i, \mu_i + 2\rho \rangle
    \bigg) \phi,
  \]
  where $r(x) = (e^x \Omega^- + \Omega^+) / (e^x - 1)$.
\end{theorem}

\section{Factorizable Systems}
\begin{remark}
  The equations we are studying take the
  form
  \[
    \frac{\partial}{\partial z_i} \phi(z_1, \dots, z_N)
    = A_i(z_1, \dots, z_N) \phi(z_1, \dots, z_N), \quad \phi \in V,\, A_i \in \End(V),
  \]
  where $\partial_i A_j \partial_j A_i + [A_i, A_j] = 0$ and
  the $A_i$ are holomorphic in
  $D \subseteq \C^N$.
\end{remark}

\begin{definition}
  We say that such a system is
  \emph{factorizable} if
  \begin{enumerate}
    \item $V = V_1 \otimes \cdots \otimes V_N$,
      where the $\{V_i\}$ are $\g$-modules,
    \item there exists a function
      $r(z) \in \g \otimes \g$ of one
      complex variable meromorphic
      in a neighborhood of the origin,
      and $\kappa \in \C$ such that
      \[
        A_i(z_1, \dots, z_N)
        = \frac{1}{\kappa}
        \bigg(
          \sum_{j > i} r_{i, j}(z_i - z_j)
          - \sum_{j < i} r_{j, i}(z_j - z_i)
          + s_i
        \bigg)
      \]
      for $s \in \g$ such that
      $[s \otimes 1 + 1 \otimes s, r(z)] = 0$.
  \end{enumerate}
\end{definition}

\begin{prop}\label{prop:classical-yang-baxter-equation}
  A factorizable system is consistent
  if
  \[
    [r_{i, j}(z_i - z_j), r_{j, k}(z_j - z_k)]
    + [r_{i, k}(z_i - z_k), r_{j, k}(z_j - z_k)]
    + [r_{i, j}(z_i - z_j), r_{i, k}(z_i - z_k)] = 0.
  \]
  These are called the
  \emph{classical Yang-Baxter equations}.
\end{prop}

\begin{remark}
  There are also
  \emph{quantum Yang-Baxter equations}, which
  is of the form
  \[
    R_{i, j}(z_i - z_j)
    = 1 + \hbar r(z_i - z_j) + O(\hbar^2),
  \]
  and one can recover the
  classical equation from it.
\end{remark}

\begin{remark}
  We can naturally take $V_i \cong U(\g)$
  to make a factorizable system, and in
  this case the
  condition in Proposition
  \ref{prop:classical-yang-baxter-equation} is
  satisfied in $U(\g)^{\otimes 3}$.
\end{remark}

\begin{definition}
  \emph{Classical $r$-matrices} are
  solutions of the classical Yang-Baxter
  equations which are meromorphic in a
  neighborhood of the origin
  with at least one point
  $z$ where $r(z)$ is nondegenerate.
\end{definition}

\begin{definition}
  Two classical $r$-matrices
  $r(z)$, $\widetilde{r}(z)$ are
  \emph{equivalent} if there is
  $\phi(x) \in \Aut(\g)$ such that
  \[
    \widehat{r}(z_1 - z_2)
    = (\phi(z_1) \otimes \phi(z_2)) r(z_1 - z_2).
  \]
\end{definition}

\begin{theorem}
  A classical $r$-matrix $r(z)$ extends
  to a meromorphic function on the
  complex plane satisfying
  $r_{1, 2}(-z) = -r_{2, 1}(z)$. Moreover,
  they fall into one of the following
  classes:
  \begin{enumerate}
    \item \emph{Rational:}
      $r(z) = c \Omega / z + P(z)$ for
      some $P \in \g \otimes \g \otimes \C[z]$
      and $c \in \C^\times$.
    \item \emph{Hyperbolic (or trigonometric):}
      $r = f(e^{az})$, where
      \[
        r(z)
        = \frac{P(e^{az})}{e^{a(n + 1)z} - 1}
        + e^{-anz} Q(e^{az})
      \]
      where $P, Q$ are polynomials
      and $r(z + 2\pi i / a) = r(z)$.
    \item \emph{Elliptic:}
      $r(z + n \omega_1) = r(z + n\omega_2) = r(z)$.
      This case exists only for
      Type A Lie algebras.
  \end{enumerate}
\end{theorem}
